\newif\ifmiddlewaresubmission
\ifdefined\MiddlewareSubmissionMode
  \middlewaresubmissiontrue
\else
  \middlewaresubmissionfalse
\fi

\newif\ifmiddlewarenamed
\ifdefined\MiddlewareNamedMode
  \middlewarenamedtrue
\else
  \middlewarenamedfalse
\fi

\newif\ifacmformat
\IfFileExists{acmart.cls}{
  \acmformattrue
  \ifmiddlewaresubmission
    \ifmiddlewarenamed
      \documentclass[sigconf,screen]{acmart}
    \else
      \documentclass[sigconf,anonymous]{acmart}
    \fi
  \else
    \documentclass[sigconf,screen]{acmart}
  \fi
}{
  \acmformatfalse
  \documentclass[11pt,a4paper]{article}
}

\ifacmformat
  \setcopyright{none}
  \ifmiddlewaresubmission
    \settopmatter{printacmref=false}
  \else
    \settopmatter{printacmref=true}
  \fi
  \renewcommand\footnotetextcopyrightpermission[1]{}
  \pagestyle{plain}
  \ifmiddlewaresubmission
    \acmConference[Middleware 2026 Submission]{ACM/IFIP/USENIX International Middleware Conference}{2026}{}
    \acmBooktitle{Middleware 2026 Submission}
  \else
    \acmConference[Research Draft]{Research Draft}{April 2026}{}
    \acmBooktitle{Research Draft}
  \fi
  \acmDOI{}
  \acmISBN{}
\else
  \usepackage[T1]{fontenc}
  \usepackage[utf8]{inputenc}
  \usepackage{lmodern}
  \usepackage[margin=2.4cm]{geometry}
  \usepackage{parskip}
  \setlength{\parindent}{0pt}
\fi

\ifacmformat\else
\usepackage[hidelinks]{hyperref}
\fi
\usepackage{booktabs}
\usepackage{tabularx}
\usepackage{enumitem}
\usepackage{array}

\setlist[itemize]{leftmargin=1.5em}
\setlist[enumerate]{leftmargin=1.7em}
\newcolumntype{Y}{>{\raggedright\arraybackslash}X}

\newcommand{\RUPStudyCount}{4}
\newcommand{\RUPOperatorCount}{2}
\newcommand{\RUPEnvironmentCount}{2}
\newcommand{\RUPTrialCount}{3}
\newcommand{\RUPTotalTrials}{12}
\newcommand{\RUPScenarioExecutions}{48}
\newcommand{\RUPSuccessfulScenarioExecutions}{48}
\newcommand{\RUPStudyStarted}{2026-04-26T16:41:00.964771+00:00}
\newcommand{\RUPStudyCompleted}{2026-04-26T17:58:52.335780+00:00}


\title{ReconcileUnderPolicy: A Repeatable Multi-Operator Benchmark for Policy-Constrained Kubernetes Recovery}
\ifmiddlewaresubmission
\subtitle{[Experimentation and Deployment]}
\fi

\ifacmformat
  \ifmiddlewaresubmission
    \ifmiddlewarenamed
      \author{Mandar Joshi}
      \affiliation{%
        \institution{Independent Researcher}
        \city{}
        \country{India}
      }
      \email{mandarjoshi1045@gmail.com}
      \ccsdesc[500]{Software and its engineering~Cloud computing}
      \ccsdesc[300]{Computer systems organization~Cloud computing}
      \keywords{Kubernetes operators, benchmarking, recovery evaluation, reproducibility, stateful systems}
    \else
    \ccsdesc[500]{Software and its engineering~Cloud computing}
    \ccsdesc[300]{Computer systems organization~Cloud computing}
    \keywords{Kubernetes operators, benchmarking, recovery evaluation, reproducibility, stateful systems}
    \fi
  \else
  \author{Mandar Joshi}
  \affiliation{%
    \institution{Independent Researcher}
    \city{}
    \country{India}
  }
  \email{mandarjoshi1045@gmail.com}
  \ccsdesc[500]{Software and its engineering~Cloud computing}
  \ccsdesc[300]{Computer systems organization~Cloud computing}
  \keywords{Kubernetes operators, benchmarking, recovery evaluation, reproducibility, stateful systems}
  \fi
\else
  \ifmiddlewaresubmission
    \date{}
  \else
    \author{Mandar Joshi \\ \texttt{mandarjoshi1045@gmail.com}}
    \date{April 2026}
  \fi
\fi

\begin{document}

\ifacmformat
\begin{abstract}
Kubernetes operators automate deployment, scaling, and recovery for complex
stateful applications, but lightweight comparative benchmarks for recovery
behavior are still uncommon. This paper presents \emph{ReconcileUnderPolicy}, a
standalone benchmark artifact for policy-constrained recovery across
\RUPOperatorCount{} stateful operators, \RUPEnvironmentCount{} execution
environments, and four recovery scenarios. The current study evaluates the
RabbitMQ Cluster Operator and CloudNativePG on both single-node and multi-node
\texttt{kind} clusters, running \RUPTrialCount{} repeated trials per study for
\RUPScenarioExecutions{} scenario executions in total. All
\RUPSuccessfulScenarioExecutions{} scenario executions converged successfully.
Median completion times reveal substantial operator-specific differences:
RabbitMQ baseline creation required 174.93\,s and 189.58\,s on single-node and
multi-node clusters, while CloudNativePG required 38.04\,s and 35.93\,s.
Delete-and-recreate recovery showed a similar gap, with RabbitMQ at 59.50\,s
and 72.81\,s versus CloudNativePG at 32.17\,s and 33.58\,s. The benchmark also
exposes warmup-sensitive variance: CloudNativePG baseline means were
substantially higher than medians because the first trial in each environment
was much slower than subsequent runs. The result is a reproducible,
artifact-backed comparative benchmark that captures both successful recovery
coverage and scenario-level timing behavior across operators and environments.
\end{abstract}
\maketitle
\else
\maketitle
\begin{abstract}
Kubernetes operators automate deployment, scaling, and recovery for complex
stateful applications, but lightweight comparative benchmarks for recovery
behavior are still uncommon. This paper presents \emph{ReconcileUnderPolicy}, a
standalone benchmark artifact for policy-constrained recovery across
\RUPOperatorCount{} stateful operators, \RUPEnvironmentCount{} execution
environments, and four recovery scenarios. The current study evaluates the
RabbitMQ Cluster Operator and CloudNativePG on both single-node and multi-node
\texttt{kind} clusters, running \RUPTrialCount{} repeated trials per study for
\RUPScenarioExecutions{} scenario executions in total. All
\RUPSuccessfulScenarioExecutions{} scenario executions converged successfully.
Median completion times reveal substantial operator-specific differences:
RabbitMQ baseline creation required 174.93\,s and 189.58\,s on single-node and
multi-node clusters, while CloudNativePG required 38.04\,s and 35.93\,s.
Delete-and-recreate recovery showed a similar gap, with RabbitMQ at 59.50\,s
and 72.81\,s versus CloudNativePG at 32.17\,s and 33.58\,s. The benchmark also
exposes warmup-sensitive variance: CloudNativePG baseline means were
substantially higher than medians because the first trial in each environment
was much slower than subsequent runs. The result is a reproducible,
artifact-backed comparative benchmark that captures both successful recovery
coverage and scenario-level timing behavior across operators and environments.
\end{abstract}
\fi

\section{Introduction}

Kubernetes operators encode operational logic into custom resources and
controllers that provision, scale, update, and recover applications through
reconciliation loops~\cite{k8s-operator}. For stateful middleware, that control
logic determines whether replicas recover correctly after restarts, whether
policy constraints block progress cleanly, and how quickly the system converges
back to a healthy state.

Recent work has shown that operator bugs are common and that automated testing
frameworks can expose subtle controller failures~\cite{acto,sieve,watchers,issta-bugs}.
At the same time, there is still a practical gap between those research
insights and a lightweight artifact that can be cloned, executed locally, and
used to compare recovery behavior across multiple real operators. Many studies
either emphasize broad bug finding or require infrastructure and orchestration
that are heavier than what a small empirical benchmark needs.

This paper addresses that gap with \emph{ReconcileUnderPolicy}, a standalone
benchmark artifact for comparative operator-recovery evaluation. The current
study evaluates two real stateful operators, RabbitMQ Cluster Operator and
CloudNativePG~\cite{rabbitmq-quickstart,cloudnativepg-docs}, across both
single-node and multi-node \texttt{kind} environments. Each study executes the
same four scenarios: baseline cluster creation, operator restart during
scale-up, ResourceQuota-blocked scale-up, and delete-and-recreate recovery.

The contribution is both empirical and methodological:

\begin{itemize}
  \item the repository executes a shared scenario suite across multiple real
  operators and environments;
  \item it captures repeated-trial timing data rather than isolated one-off
  runs;
  \item it packages machine-readable results and LaTeX-ready tables directly
  from the benchmark outputs;
  \item it exposes where recovery behavior is stable, where it is
  operator-specific, and where warmup effects distort means.
\end{itemize}

The resulting paper is therefore not a claim about all Kubernetes operators.
Instead, it is a reproducible comparative benchmark study that demonstrates how
policy-constrained recovery can be evaluated with a small but meaningful matrix
of operators and environments.

\section{Benchmark Scope}

\subsection{Research Questions}

The benchmark focuses on three questions:

\begin{enumerate}
  \item Can a lightweight local artifact repeatedly execute non-trivial
  recovery scenarios across multiple real stateful operators?
  \item How do recovery times differ across operators and environments when the
  scenario suite is held constant?
  \item What kinds of run-to-run variability become visible once repeated
  trials are collected rather than single-run anecdotes?
\end{enumerate}

\subsection{Study Matrix}

The present evaluation spans \RUPStudyCount{} studies:

\begin{itemize}
  \item operators: RabbitMQ Cluster Operator and CloudNativePG;
  \item environments: single-node \texttt{kind} and multi-node \texttt{kind};
  \item trials: \RUPTrialCount{} full trials per operator/environment pair;
  \item scenarios: four recovery scenarios per trial.
\end{itemize}

This produces \RUPTotalTrials{} trial runs and \RUPScenarioExecutions{} total
scenario executions.

\subsection{Scenario Suite}

Each trial executes four scenarios in sequence:

\begin{enumerate}
  \item \textbf{baseline\_create}: create a one-replica managed cluster and
  measure time to first healthy reconciliation;
  \item \textbf{restart\_during\_scale}: scale from one to two replicas while
  deleting the operator pod mid-reconcile;
  \item \textbf{quota\_blocked\_scale}: apply a temporary ResourceQuota that
  blocks creation of the next pod during scale-up, then remove it and measure
  recovery;
  \item \textbf{delete\_and\_recreate}: delete the managed cluster, wait for
  cleanup, recreate it, and measure time back to a healthy one-replica state.
\end{enumerate}

\section{Artifact Design}

\subsection{Repository Structure}

The repository is organized into three main areas:

\begin{itemize}
  \item \texttt{vendor/cluster-operator-main}: vendored RabbitMQ Cluster
  Operator source used to render local install manifests;
  \item \texttt{vendor/cloudnative-pg}: vendored CloudNativePG release assets;
  \item \texttt{benchmark/reconcileunderpolicy/}: benchmark scripts,
  environment definitions, study outputs, and reproduction steps;
  \item \texttt{paper/}: LaTeX source, generated tables, and bibliography.
\end{itemize}

\subsection{Execution Environment}

The benchmark runs on two dedicated \texttt{kind} profiles: a single-node
cluster and a three-node cluster with one control-plane node and two workers.
RabbitMQ and CloudNativePG are both exercised with the same per-instance
resource envelope: 200m CPU and 512Mi memory requests, plus 500m CPU and 1Gi
memory limits. This keeps the scenarios comparable while remaining feasible on
commodity local hardware.

RabbitMQ installation is rendered locally from the vendored source tree.
CloudNativePG installation uses the vendored \texttt{v1.29.0} release manifest
applied with server-side \texttt{kubectl}. Managed workload images are
\texttt{rabbitmq:4.1.3-management} for RabbitMQ. CloudNativePG uses
\texttt{ghcr.io/cloudnative-pg/postgresql:17}.

\subsection{Outputs}

Each study run creates a timestamped directory containing:

\begin{itemize}
  \item per-trial summaries,
  \item per-scenario snapshots of pods, PVCs, services, and events,
  \item aggregate JSON, CSV, and Markdown summaries,
  \item generated LaTeX include files for the manuscript tables.
\end{itemize}

\section{Methodology}

Each study begins from a clean workload state, ensures the operator deployment
is healthy, and then runs the four scenarios in order. Success is defined as
convergence to the expected healthy replica count for the scenario under test.

For each scenario, the benchmark records:

\begin{itemize}
  \item end-to-end completion time,
  \item final ready replica count,
  \item scenario-specific metadata such as the deleted operator pod name or the
  ready-replica count observed while the quota block was active.
\end{itemize}

Aggregated statistics are computed per scenario for each operator/environment
pair:

\begin{itemize}
  \item successful trials out of total trials,
  \item mean completion time,
  \item median completion time,
  \item standard deviation,
  \item minimum and maximum completion times.
\end{itemize}

Both mean and median are reported because several scenarios exhibit clear
first-trial skew. Means capture the overall cost actually observed in the
study, while medians better reflect steady repeated behavior when a single cold
start dominates one run.

\section{Results}

Table~\ref{tab:median-results} summarizes the per-study median completion times
for the four scenarios. Table~\ref{tab:detailed-results} then reports the full
aggregate statistics for each operator/environment pair.

\begin{table}[ht]
\centering
\small
\caption{Median completion time in seconds by operator and environment}
\label{tab:median-results}
\begin{tabular}{llrrrr}
\toprule
\textbf{Operator} & \textbf{Env.} & \textbf{Baseline} & \textbf{Restart} & \textbf{Quota} & \textbf{Delete} \\
\midrule
RabbitMQ & Single-node & 174.93 & 77.61 & 79.54 & 59.50 \\
RabbitMQ & Multi-node & 189.58 & 88.74 & 92.97 & 72.81 \\
CloudNativePG & Single-node & 38.04 & 27.72 & 54.27 & 32.17 \\
CloudNativePG & Multi-node & 35.93 & 31.50 & 54.13 & 33.58 \\
\bottomrule
\end{tabular}
\end{table}

\begin{table*}[t]
\centering
\scriptsize
\caption{Aggregate statistics across \RUPTrialCount{} trials per study}
\label{tab:detailed-results}
\begin{tabularx}{\textwidth}{llYcccc}
\toprule
\textbf{Operator} & \textbf{Env.} & \textbf{Scenario} & \textbf{Success} & \textbf{Mean (s)} & \textbf{Stddev (s)} & \textbf{Range (s)} \\
\midrule
RabbitMQ & Single-node & Baseline create & 3/3 & 175.00 & 0.20 & 174.80--175.27 \\
RabbitMQ & Single-node & Restart during scale & 3/3 & 77.16 & 1.48 & 75.17--78.71 \\
RabbitMQ & Single-node & Quota-blocked scale & 3/3 & 80.23 & 1.09 & 79.37--81.77 \\
RabbitMQ & Single-node & Delete and recreate & 3/3 & 59.49 & 0.07 & 59.41--59.57 \\
RabbitMQ & Multi-node & Baseline create & 3/3 & 192.33 & 3.96 & 189.49--197.93 \\
RabbitMQ & Multi-node & Restart during scale & 3/3 & 149.33 & 85.80 & 88.58--270.66 \\
RabbitMQ & Multi-node & Quota-blocked scale & 3/3 & 93.53 & 0.87 & 92.85--94.76 \\
RabbitMQ & Multi-node & Delete and recreate & 3/3 & 72.82 & 0.13 & 72.67--72.98 \\
CloudNativePG & Single-node & Baseline create & 3/3 & 103.39 & 94.31 & 35.37--236.76 \\
CloudNativePG & Single-node & Restart during scale & 3/3 & 28.78 & 1.64 & 27.53--31.09 \\
CloudNativePG & Single-node & Quota-blocked scale & 3/3 & 52.93 & 2.07 & 50.00--54.51 \\
CloudNativePG & Single-node & Delete and recreate & 3/3 & 33.05 & 1.25 & 32.16--34.81 \\
CloudNativePG & Multi-node & Baseline create & 3/3 & 114.38 & 111.13 & 35.68--271.54 \\
CloudNativePG & Multi-node & Restart during scale & 3/3 & 102.06 & 101.44 & 29.16--245.51 \\
CloudNativePG & Multi-node & Quota-blocked scale & 3/3 & 54.71 & 1.02 & 53.85--56.15 \\
CloudNativePG & Multi-node & Delete and recreate & 3/3 & 32.94 & 0.92 & 31.63--33.60 \\
\bottomrule
\end{tabularx}
\end{table*}

Four observations stand out.

\textbf{All scenario executions converged successfully.}
Every one of the \RUPScenarioExecutions{} scenario executions succeeded,
including operator restart during scale-up, temporary policy blocking, and full
delete-and-recreate recovery. The artifact therefore exercises disruptive
control-plane events without requiring manual intervention or operator-specific
special casing in the paper workflow.

\textbf{CloudNativePG recovered faster than RabbitMQ in baseline and
delete-and-recreate scenarios.}
On single-node \texttt{kind}, median baseline creation was 38.04\,s for
CloudNativePG versus 174.93\,s for RabbitMQ. On multi-node \texttt{kind}, the
same comparison was 35.93\,s versus 189.58\,s. Delete-and-recreate showed a
similar pattern: 32.17\,s and 33.58\,s for CloudNativePG versus 59.50\,s and
72.81\,s for RabbitMQ. That gap suggests the benchmark is measuring genuine
operator-specific behavior rather than only generic cluster overhead.

\textbf{Environment effects were operator-dependent rather than uniform.}
RabbitMQ slowed down on the multi-node profile across all four scenarios,
including restart-during-scale median growth from 77.61\,s to 88.74\,s and
delete-and-recreate growth from 59.50\,s to 72.81\,s. CloudNativePG was much
less sensitive in median terms: baseline medians were 38.04\,s and 35.93\,s,
while delete-and-recreate medians were 32.17\,s and 33.58\,s. The environment
tax is therefore not a fixed property of the benchmark harness; it depends on
the operator and scenario.

\textbf{Means exposed warmup-sensitive outliers that medians would hide.}
CloudNativePG baseline creation had means of 103.39\,s and 114.38\,s but
medians of only 38.04\,s and 35.93\,s, driven by a much slower first trial in
each environment. RabbitMQ multi-node restart during scale showed a similar
pattern, with an 88.74\,s median but a 149.33\,s mean due to one 270.66\,s
trial. Reporting both measures is therefore important: the means preserve the
full observed cost, while the medians expose the more typical steady-state
recovery time.

\textbf{Quota-blocked recovery was consistently observable and stable.}
All four studies successfully stalled progress while the ResourceQuota was in
place and then recovered once the quota was removed. Standard deviations stayed
low for this scenario across the matrix, ranging from 0.87\,s to 2.07\,s. This
makes quota-blocked scale-up a useful cross-operator stressor for future work
because it creates a real controller-visible perturbation without modifying the
operators themselves.

\section{Discussion}

\subsection{What the Study Establishes}

This work does not claim to characterize the entire operator ecosystem. It does
establish that a lightweight repository can support a non-trivial comparative
matrix over real operators and environments while still remaining executable on
local infrastructure. The benchmark now moves beyond a single-operator pilot:
it produces comparative tables, captures repeated-trial variance, and exposes
where operator recovery behavior diverges under the same perturbation suite.

\subsection{Threats to Validity}

The claims of the paper are intentionally scoped, and the main validity limits
concern generalization.

\begin{itemize}
  \item \textbf{External validity.} The study covers only two operators, so the
  measured recovery behaviors should not be generalized to all operator
  ecosystems without broader sampling.
  \item \textbf{Environment validity.} Both environments are local
  \texttt{kind} clusters rather than production deployments, so absolute
  timings are environment-specific even though the workflows are realistic.
  \item \textbf{Policy validity.} The current perturbation suite centers on
  ResourceQuota-based interference and controller restarts. Other policy
  mechanisms such as admission rejection or RBAC denial are not yet included.
  \item \textbf{Construct validity.} The benchmark measures convergence and
  recovery timing, not deeper workload semantics beyond operator health and
  replica readiness.
  \item \textbf{Warmup effects.} Several scenarios showed first-trial skew,
  which means future studies should control for image pull, initialization, or
  storage warmup when comparing absolute means.
\end{itemize}

These constraints bound the conclusions, but they do not undermine the central
contribution: a reproducible comparative benchmark for policy-constrained
operator recovery with real execution traces and repeated measurements.

\subsection{Next Steps}

The most promising extensions are:

\begin{enumerate}
  \item add more operators while preserving the same scenario suite;
  \item incorporate additional policy mechanisms such as admission denial or
  RBAC restriction;
  \item extend the environments beyond local \texttt{kind} clusters;
  \item correlate timing differences with richer controller-log and event-trace
  analysis.
\end{enumerate}

\section{Related Work}

Acto~\cite{acto} introduced automated end-to-end testing for operation
correctness in cloud system management. Sieve~\cite{sieve} extended this line
of work with perturbation-based testing for Kubernetes controllers. The NSDI
2026 reliability study~\cite{watchers} and the ISSTA 2024 operator bug
study~\cite{issta-bugs} provide broader motivation by showing that operator
failures are both common and consequential.

The present work is narrower than those systems papers. Its contribution lies
in packaging a runnable comparative benchmark artifact around real stateful
operators, exposing policy-constrained recovery behavior across repeated trials,
and turning those outputs directly into paper-ready comparative tables.

\section{Conclusion}

\emph{ReconcileUnderPolicy} now provides a repeatable comparative benchmark for
policy-constrained recovery in Kubernetes operators. Across \RUPStudyCount{}
studies, \RUPTotalTrials{} repeated trial runs, and \RUPScenarioExecutions{}
scenario executions, the artifact successfully evaluated RabbitMQ Cluster
Operator and CloudNativePG on both single-node and multi-node \texttt{kind}
profiles.

The resulting data show that recovery behavior is meaningfully
operator-dependent: CloudNativePG converged substantially faster than RabbitMQ
in baseline and delete-and-recreate scenarios, RabbitMQ exhibited a clearer
multi-node penalty, and several scenarios displayed warmup-sensitive outliers
that would be missed by single-run evaluation. The repository therefore
supports a stronger publication claim than a pilot artifact alone: it offers a
reproducible methodology, real comparative data, and a reusable baseline for
larger studies of policy-constrained operator recovery.

\ifacmformat
\bibliographystyle{ACM-Reference-Format}
\else
\bibliographystyle{plain}
\fi
\bibliography{references}

\end{document}
